% Generated by `taskgen course examples` from dlp.rho_distinguished@sem2 (seed 11).
% Do not edit: change examples.yaml and regenerate.
\begin{example}[выделенные точки]\label{ex:rho-distinguished}
Решить задачу дискретного логарифмирования $64^x \equiv 36\ (\textrm{mod}\ 431)$, $\ord(64) = 43$, $\rho$-методом с выделенными точками. Два вычислителя параллельно (по одному шагу за такт) строят последовательности $x_{t+1} = f(x_t)$ и $y_{t+1} = f(y_t)$, где\[f(x) = \begin{cases} x^2, & x \equiv 0,\\ 36\, x, & x \equiv 1,\\ 64\, x, & x \equiv 2 \end{cases} \pmod 3, \qquad (\alpha, \beta) \mapsto \begin{cases} (2\alpha, 2\beta),\\ (\alpha, \beta + 1),\\ (\alpha + 1, \beta), \end{cases}\]с показателями по модулю $43$; $x_0 = 64^{40} \cdot 36^{29} \equiv 384$, $y_0 = 64^{3} \cdot 36^{21} \equiv 16$. Точка \emph{выделенная}, если её 3 младших двоичных разряда нулевые, т.\,е. она делится на $8$. Выделенные точки (с показателями) отправляются на сервер; работа прекращается, как только какая-то выделенная точка придёт на сервер повторно. Выписать таблицу вычислений и найти $x$.

\par\noindent\textit{Решение.} Таблица (выделенные точки подчёркнуты; показатели $x_t = g^{\alpha_t} h^{\beta_t}$, $y_t = g^{\gamma_t} h^{\delta_t}$):\[\begin{array}{r|rrr|rrr}t & x_t & \alpha_t & \beta_t & y_t & \gamma_t & \delta_t \\ \hline 0 & \underline{384} & 40 & 29 & \underline{16} & 3 & 21 \\ 1 & 54 & 37 & 15 & 145 & 3 & 22 \\ 2 & 330 & 31 & 30 & \underline{48} & 3 & 23 \\ 3 & \underline{288} & 19 & 17 & 149 & 6 & 3 \\ 4 & \underline{192} & 38 & 34 & 54 & 7 & 3 \\ 5 & 229 & 33 & 25 & 330 & 14 & 6 \\ 6 & 55 & 33 & 26 & \underline{288} & 28 & 12 \\ \end{array}\]Содержимое сервера до совпадения: $384$ ($x_{0}$: $\alpha = 40$, $\beta = 29$); $16$ ($y_{0}$: $\gamma = 3$, $\delta = 21$); $48$ ($y_{2}$: $\gamma = 3$, $\delta = 23$); $288$ ($x_{3}$: $\alpha = 19$, $\beta = 17$); $192$ ($x_{4}$: $\alpha = 38$, $\beta = 34$). При $t = 6$ точка $288$ приходит повторно (от второго вычислителя: последовательности слились).\par Из $g^{19} h^{17} = g^{28} h^{12}$ и $h = g^x$ получаем $(17 - 12)\, x \equiv 28 - 19\ (\textrm{mod}\ 43)$, т.\,е. $5\, x \equiv 9\ (\textrm{mod}\ 43)$. $\gcd(5, 43) = 1$, решение единственно: $x \equiv 9 \cdot 5^{-1}$.  Обратный к $5$ по модулю $43$ (расширенный алгоритм Евклида для $(43, 5)$):
\[\begin{array}{r|r|r|r}
r_i & q_i & u_i & v_i \\ \hline
43 &  & 1 & 0 \\
5 &  & 0 & 1 \\
3 & 8 & 1 & -8 \\
2 & 1 & -1 & 9 \\
1 & 1 & 2 & -17 \\
0 & 2 & -5 & 43 \\
\end{array}\qquad r_i = u_i \cdot 43 + v_i \cdot 5\]

$5^{-1} \equiv -17 \equiv 26\ (\textrm{mod}\ 43)$.
\par $\tilde x \equiv 9 \cdot 26 \equiv 19\ (\textrm{mod}\ 43)$. Проверка: $64^{19} \equiv 36\ (\textrm{mod}\ 431)$.\[x = 19.\]
\end{example}
